\section{Lattice composition with matrix symbols}\label{app:libcalculus}

The Weyl calculus is classical; see \cite{HSII,RammalBellissard}.
We record the elementary discrete calculation to specify matrix
order, shifted domains and the errors used here. For
$a(k,t)=\sum_m a_m(t)e^{imk}$, put $t_j=t^\nu+j\delta$ and
\[
 (\operatorname{Op}(a)v)_j
       =\sum_m a_m(t_j-m\delta/2)v_{j-m},\qquad
 N_r^p(a)=\sum_m(1+|m|)^p\sup_I\|\partial_t^r a_m\|.
\]
In estimates requiring an off-cell factor, omit the zero Fourier
mode. The interval $I$ contains every shifted segment below.
Absolute row and column sums give $\|\operatorname{Op}(a)\|\le N_0^0(a)$.
For the symmetric endpoint-average quantization,
\begin{equation}\label{eq:libendpoint}
 \|\operatorname{Op}^{\rm sym}(a)-\operatorname{Op}(a)\|
       \le\delta^2 N_2^2(a)/8.
\end{equation}

\begin{lemma}[Exact product and second-order remainder]
\label{lem:libcalculus}
Matrix order is preserved in
\begin{equation}\label{eq:libtwist}
 (a\mathbin\# b)_\ell(t)
   =\sum_{m+n=\ell}a_m(t+n\delta/2)b_n(t-m\delta/2),\qquad
 \operatorname{Op}(a)\operatorname{Op}(b)
       =\operatorname{Op}(a\mathbin\# b).
\end{equation}
With $\{a,b\}=a_kb_t-a_tb_k$,
\begin{equation}\label{eq:libTaylorbound}
 \left\|\operatorname{Op}(a)\operatorname{Op}(b)
       -\operatorname{Op}\left(ab+\frac{i\delta}2\{a,b\}\right)\right\|
 \le\frac{\delta^2}8\bigl[
 N_2^0(a)N_0^2(b)+2N_1^1(a)N_1^1(b)+N_0^2(a)N_2^0(b)\bigr].
\end{equation}
If $P_J$ is the cell projection onto a finite well, then
\begin{equation}\label{eq:libinterface}
 (P_J\operatorname{Op}(a)P_J)(P_J\operatorname{Op}(b)P_J)
 =P_J\operatorname{Op}(a\mathbin\# b)P_J
  -P_J\operatorname{Op}(a)(1-P_J)\operatorname{Op}(b)P_J.
\end{equation}
The last term uses the full-lattice factors before compression.
\end{lemma}

\begin{proof}
Multiplying kernels at $(j,j-m,j-m-n)$ gives
\eqref{eq:libtwist}. Apply Taylor's formula in $u\in[0,1]$ to
$a_m(t+un\delta/2)b_n(t-um\delta/2)$.
The first derivative gives $(i\delta/2)\{a,b\}$.
The three second derivatives have coefficients $n^2,-2mn,m^2$.
Integrating their norms and summing both absolute Schur sums proves
\eqref{eq:libTaylorbound}. This uses only second derivatives and
covers negative $\delta$. Equation~\eqref{eq:libendpoint} is the
trapezoid-midpoint bound. Insert $1=P_J+(1-P_J)$ between full
factors for~\eqref{eq:libinterface}.
\end{proof}

Here is the family-specific control of these seminorms.
With the two finite anchor lifts at shorter-arc distance $r\le q/2$
and $a_0=L-4\tau$,
\begin{equation}\label{eq:libconvolution}
 \sum_x e^{-a_0\dist(x,\mathcal A)}
        e^{-a_0\dist(x,\mathcal A+mq)}
 \le C(|m|q+1)e^{-a_0(|m|q-r)},\qquad m\ne0.
\end{equation}
Fixed differentiated versions follow from~\eqref{eq:libtails}.
The matrix coefficients of $h$ are inner products of two frame
vectors with the nearest-neighbor frozen operator. The orbital
matrix also has one position factor, costing $q(1+|m|)$.
Thus, for the finite $p,r'$ used in the preceding estimates,
\begin{align}
 N_{r'}^p(Y-Y_0^{\rm Fourier})
       &\le z_\gamma\mathcal R_qe^{-a_0(q-r)},\\
 N_{r'}^p(W-W_0^{\rm Fourier})
       &\le qz_\gamma\mathcal R_qe^{-a_0(q-r)},\\
 N_{r'}^0(Y)&\le z_\gamma\mathcal R_q,\qquad
 N_{r'}^0(W)\le qz_\gamma\mathcal R_q.\label{eq:libseminorms}
\end{align}
The Fourier subscripts in this display are unrelated to the energy
parameter $Y_0$ or the floor $W_0$.
The same common periodic gauge is used for every coefficient.
The proof is~\eqref{eq:libconvolution}, followed by summation over
$m$; finite extra polynomial factors are absorbed in $\mathcal R_q$.
It relies on the localized frame, not merely on $Z^2=xI_2$.

For Fourier truncation at $|m|\le q$, enlarge the phase interval by
$[-3q|\delta|/2,3q|\delta|/2]$. It then contains all product
midpoints and shifted segments. This enlargement is smaller than
the reserved frame-chart margin for $b\ge2\xi$.
The omitted Fourier tail has its absolute weighted sum retained.
Every nonzero term crossing a shell or well-end interface contains
two off-cell factors in~\eqref{eq:libinterface}. Its norm is at most
$C\overline H(\overline H+h_\vartheta)$.
It is not part of an $O(\delta^2)$ Taylor remainder.

For completeness, put $u=Y_0-\vartheta$, so the two symbols in
\eqref{eq:libproduct} are $b=\widetilde Y+\varepsilon W+u$ and
$a=\widetilde Y+\varepsilon W-u$.
Then $ba=\Phi I_2+\varepsilon\{\widetilde Y,W\}_+
+\varepsilon^2W^2$ and
$\{b,a\}=[y_k,y_t]+2\{y,\vartheta\}$.
Equations~\eqref{eq:libTaylorbound} and~\eqref{eq:libseminorms}
prove the product and its $\kappa\overline H\mathcal R_q$ error.
This also identifies the anti-Hermitian term that must be retained.

\section{The localized three-block boundary calculation}
\label{app:libboundary}

We give the end correction separately because the global estimate
$c_{\rm end}^2/m\ll A$ is not required by the argument.
Fix a fine window and use its shell penalty and capped weight.
Write $K=e^\varphi(H+M_EP_{\mathcal S}-E)e^{-\varphi}$.
Split the frame projection as $\Pi=\Pi_I+\Pi_B$, where $B$ consists
of end strips beyond the deepest shell and its padding, and
$I$ is the remaining frame. Keep Fourier and interwell tails as
errors in their assigned budgets. The weight is constant at this
end seam. Write $F=K_{\Pi\Pi}$.

The shell gap, compression estimate and weighted frame estimate give
$\|F^{-1}\|\le1/(2A)$, with the same adjoint statement.
At an end cell, $U_{\rm end}\asymp\Lambda e^{-\xi q/2}$ and the
spectral depth is $e\le4\Lambda e^{-\xi q}$.
The product lower bound and its partner normalization give
\begin{equation}\label{eq:libendmargin}
 \begin{split}
 d:=\sigma_{\min}(F_{BB})
 &\ge\frac{c(U_{\rm end}-e)}
 {z_\gamma(\sqrt{\mathrm{ex}+U_{\rm end}}
                         +\sqrt{\mathrm{ex}+e})}\\
 &\ge\frac{cU_{\rm end}}{z_\gamma\sqrt{\mathrm{ex}+U_{\rm end}}}
 \ge e^{-(\xi/2+\text{chosen loss})q}.
 \end{split}
\end{equation}
When $\mathrm{ex}$ dominates, replacing this divided difference by
$\sqrt{U_{\rm end}}/z_\gamma$ would be incorrect.
The penalty is zero on $B$.
Also
\[
 \|F_{IB}\|,\|F_{BI}\|\le t_{\rm end}
       =\overline H/z_\gamma\le e^{-(L/2-2\tau)q}\mathcal R_q,
 \qquad \|K_{QQ}^{-1}\|\le2/m.
\]
Let $c_0$ now denote the complete weighted bulk cross-block bound,
so $c_0^2/m\le A/50$ after assigning the budgets and
$c_0=o(m)$. The end cross blocks are bounded by $c_0+c_{\rm end}$.
The weighted $c_0$ need not be $O(|\delta|\mathcal R_q)$.

Eliminate $Q$, writing $R=K_{QQ}^{-1}$ and
$S=F-G$, $G=K_{\Pi Q}RK_{Q\Pi}$.
Then
\[
 \|G_{II}\|\le2c_0^2/m,\quad
 \|G_{IB}\|,\|G_{BI}\|\le2c_0(c_0+c_{\rm end})/m,\quad
 \|G_{BB}\|\le2(c_0+c_{\rm end})^2/m.
\]
Set
\[
 r_1=(c_0+c_{\rm end})^2/(md),\quad
 r_2=t_{\rm end}/d,\quad
 r_3=c_{\rm end}^2t_{\rm end}^2/(md^2A).
\]
One has $A/d=o(1)$, by~\eqref{eq:libunit}, the raw compression
errors and~\eqref{eq:libendmargin}. Hence the bulk part of $r_1$ is
at most $A/(50d)=o(1)$. Its end part tends to zero.
Furthermore
\begin{equation}\label{eq:librthree}
 r_2=o(1),\qquad
 r_3\le e^{(-L+3\xi/2+4\tau)q}\mathcal R_q^C=o(1).
\end{equation}
The last bound uses $A\ge\epsilon_{\rm ref}
\ge e^{-(L+\text{chosen loss})q}$ and $\xi$ sufficiently small.
All estimates hold for adjoints too.

For $q$ large, $S_{BB}$ is invertible with norm of its inverse at
most $2/d$. The resolvent identity gives
\[
 S_{BB}^{-1}-F_{BB}^{-1}=F_{BB}^{-1}G_{BB}S_{BB}^{-1}.
\]
The difference of interior Schur complements is exactly
\begin{equation}\label{eq:libSchurdifference}
\begin{split}
 S/S_{BB}-F/F_{BB}={}&-G_{II}
 -F_{IB}(S_{BB}^{-1}-F_{BB}^{-1})F_{BI}\\
 &+G_{IB}S_{BB}^{-1}F_{BI}
   +F_{IB}S_{BB}^{-1}G_{BI}
   -G_{IB}S_{BB}^{-1}G_{BI}.
\end{split}
\end{equation}
Both mixed $I$--$Q$--$B$ paths are present.
The second term costs at most
$8r_2^2c_0^2/m+8r_3A$.
The two mixed terms together cost at most
\[
 8r_2c_0^2/m+
 8\sqrt{c_0^2/m}\sqrt{r_3A}.
\]
The last term costs at most $8r_1c_0^2/m$.
Together with $\|G_{II}\|\le A/25$, their sum is at most $A/10$
for all large $q$.

The inverse of the unperturbed interior Schur complement is the
$I,I$ block of $F^{-1}$. It therefore has norm at most $1/(2A)$;
no positivity of $F$ is assumed. Thus the perturbed interior inverse
has norm at most $1/(1.9A)$.
Since $r_2=o(1)$ and $A/d=o(1)$, block reconstruction gives
$\|S^{-1}\|\le0.6/A$ eventually. The strict little-$o$ comparison
$A/d=o(1)$ is needed for this constant; $A\le d/4$ alone would not
suffice. Finally $A=o(m)$ and both full cross blocks are $o(m)$,
so reconstructing $Q$ gives $\|K^{-1}\|\le1/A$.
The same calculation for $K^*$ supplies the two-sided inverse
statement. A small complex strip, for example
$|\operatorname{Im}z|\le A/4$ after a fixed contraction of the real
fine window, has inverse $O(1/A)$ by the resolvent identity.

\section{From weighted gaps to launch-shell confinement}
\label{app:liblocal}

We supply the functional-calculus details used in
Proposition~\ref{prop:FPtransport}.
The physical Green, synthesis and orthogonalization kernels are
those of \cite[Appendices B--D]{DMband}, extended to rank at most two
in Lemma~\ref{lem:libframes}. Let $U_{\mathcal S}$ synthesize shell
vectors. The weight is zero on their cells. For a tail at cell
distance $d$ from such a cell,
\[
 |e^{\pm\varphi(x)}-1|\le C\mu d e^{\mu d}.
\]
Combining this with the physical frame tails gives an absolute
kernel majorant with one tail from a shell cell to a jump bond.
Its row and column sums are $C\mu\sqrt F$.
Compressing to the frame multiplies by a second localized kernel;
convolution and both Schur sums give $C\mu F$.
Both endpoints of every jump bond are included. The number of
shell components is bounded, so the constants are independent of
their lengths. Across wells retain the physical interwell tail
multiplied by $e^{\mu D_{\rm cap}}$, as in
\cite[Appendix D, Boundary leakage with the correct rate]{DMband}.
It remains negligible. This proves~\eqref{eq:libonetails}, including
both weight signs. Expanding
$e^\varphi P_{\mathcal S}e^{-\varphi}-P_{\mathcal S}$ gives the
one-tail and two-tail penalty bounds used in
Proposition~\ref{prop:libweighted}.

The same proof in rank one gives the top, bottom and W estimates.
At low $b$ the normalized compression errors can grow.
Use $A=D/|P'(E_a)|$, not a dimensionless margin exceeding the
range of a previously normalized gap lemma. The coarse shell width
and the larger padding proved above still give~\eqref{eq:FPwidth}.
The raw top-frame estimates require only $b\ge2\xi$: on the IMS
patches the freezing error $C\sqrt{|\delta|/q}$ is $o(m)$.
The small-$\kappa$ Mathieu ladder and its quantization assumptions
are unused in this extension.

\subsection{Static localization and the initial tail}

Let $H_O=H+M_EP_{\mathcal S}$ and let $g$ be a smooth fine filter
supported in the contracted gap window of $H_O$.
The frame, complement and boundary calculations above give
\[
 g(H_O)=0,\qquad
 \|e^\varphi(H_O-z)^{-1}e^{-\varphi}\|\le C/A
\]
on the small complex strip used for an almost-analytic extension.
Its vertical width and all required inverse constants are
$e^{O(q)}$.
The resolvent identity and Helffer--Sj\"ostrand formula give
\begin{equation}\label{eq:libstatic}
 \|\un_{\mathcal O}g(H)\|
 \le C(g)M_EA^{-1}e^{-\mu p_{\rm pad}}
 \le\exp(-e^{\tau q})
\end{equation}
after harmless contractions and polynomial factors.
Here $\mathcal O$ is outside the padded union of shell cells.
Although the penalty is nonlocal, its weighted synthesis is bounded
by~\eqref{eq:libonetails}; this is the needed right-hand factor in
the resolvent identity.
Choose a wider filter $g_2=1$ on a neighborhood of $\supp g$, still
inside a spare gap buffer, with fixed-order derivatives $e^{O(q)}$.
Constructing the initial gap with margin $4A$ and using a narrower
fine window leaves this buffer. This avoids assuming that an
arbitrary compact support stays a fixed distance from a window end.

Merge padded shell intervals that meet or are separated by less
than a fixed multiple of the padding. Their number is bounded, so
each merged component still has width bounded by a constant times
\eqref{eq:FPwidth}. If the launch $x_0$ lies outside the padded union,
$\|g(H)\delta_{x_0}\|$ is bounded by~\eqref{eq:libstatic}.
Otherwise choose the merged component containing it and a position
cutoff $\chi$ whose boundary bonds lie in $\mathcal O$.
In the functional-calculus formula for $[g(H),\chi]$, split
\[
 (H-z)^{-1}=(H-z)^{-1}g_2(H)+(H-z)^{-1}(1-g_2(H)).
\]
The term with both resolvents in the second part is analytic near
$\supp g$ and integrates to zero. Each remaining term contains
$[H,\chi]g_2(H)$ or its adjoint, which is bounded by
\eqref{eq:libstatic}. The functional-calculus factors are
$e^{O(q)}$, so the initial mass outside this component is doubly
exponentially small.

For $\psi_t=e^{-itH}g(H)\delta_{x_0}$,
\[
 \left|\frac{d}{dt}\langle\psi_t,(1-\chi)\psi_t\rangle\right|
 \le\|[H,\chi]g_2(H)\|\,\|\psi_t\|^2.
\]
Thus the outside mass is at most
$\exp(-e^{c\xi q})(1+|t|)$ for a fixed $c>0$, after reducing the
constant in the exponent. This is confinement about the launch
shell. Its distance from the well bottom is not included in its
width. Summing the $e^{O(q)}$ fine windows preserves every fixed
inverse-power tail on $|t|\le e^{Bq}$.

\subsection{The interior functional-calculus bridge}

Let $f$ be a smooth filter supported in the interior class of
\eqref{eq:libinterior}, with derivatives $e^{O(q)}$.
Insert a smooth spectral-neighborhood cutoff equal to one on the
spectrum. For $|t|\le e^{Bq}$, all fixed derivatives of
$h_t(E)=e^{-itE}f(E)$ are $e^{O(q)}$.
In an almost-analytic formula for $h_t(H)$ choose fixed extension
order $r$ and vertical cutoff $y_0=e^{-dq}$.
Fix the desired error $e^{-Mq}$ and choose a constant $B'$ large.
For $e^{-B'q}\le|\operatorname{Im}z|\le y_0$, use the interior
resolvent estimate at auxiliary time $1/|\operatorname{Im}z|$.
Its doubly exponential error absorbs every $e^{O(q)}$ factor.
For $|\operatorname{Im}z|<e^{-B'q}$, the plain resolvent bound and
$\bar\partial\widetilde h_t=O(e^{C_rq}|\operatorname{Im}z|^r)$
give at most $Ce^{(C_r-rB')q}$.
Choose $B'$ so that this is below $e^{-Mq}$.
The argument treats the negative half-plane by conjugation.
It proves the interior filtered-packet estimate on every fixed
exponential time window. It asserts no longer time range.
This completes the analytic passage required by
Proposition~\ref{prop:FPtransport}.
